\section{Реализация на основните подсистеми}
\label{sec:core_implementation}

След като архитектурният модел е дефиниран, следващият въпрос е как той се реализира на ниво конкретни C++ конструкции и как отделните подсистеми си взаимодействат. Настоящата глава разглежда основните building blocks на библиотеката и ги проследява от entity описанията до изпълнението на заявките, подготвените заявки, миграциите и проверката на ограниченията.

Тук фокусът вече не е върху слоевете като архитектурна схема, а върху тяхната реална механика. Как точно се натрупват \code{where} клаузите? Как query обектът се запазва за многократна употреба? Как \code{migrate<DB>} превръща разлика между entity модел и \code{live\_schema} в DDL операции? Как нишковият слой гарантира коректна употреба на връзките? Отговорът на тези въпроси показва доколко библиотеката е реално завършена като инженерна система.

\subsection{Repository-backed mini case study: \code{User} и \code{Post}}

За да остане изложението обвързано с реалния код, тази глава използва като опорна mini case study типовете \code{User} и \code{Post} от \path{tests/integration/test_mockdb.cpp}. Те са особено полезни, защото съчетават скаларни полета, table naming и join сценарий. Същият пример по-късно се използва и за prepared query, result typing и SQL rendering assertions.

\begin{lstlisting}[language=C++,caption={Entity типове от \protect\path{tests/integration/test_mockdb.cpp}},label={lst:mockdb_entities}]
struct Post
{
    orm::property<int, "id">             id;
    orm::property<int, "author_id">      author_id;
    orm::property<std::u8string, "body"> body;
};

struct User
{
    orm::property<int, "id">            id;
    orm::property<std::u8string, "name"> name;
    orm::property<double, "score">      score;
};
\end{lstlisting}

Listing~\ref{lst:mockdb_entities} показва не само как се описва entity моделът, но и как той непосредствено участва в тестовете. \code{table\_name\_trait<User>} и \code{table\_name\_trait<Post>} фиксират логическите контейнери \code{"users"} и \code{"posts"}, а member pointers като \code{\&User::id} и \code{\&Post::author\_id} по-късно участват в query builder-а. Това е добра илюстрация на един от централните мотиви на тезата: моделът, query формата и execution тестовете използват едни и същи compile-time артефакти.

\subsection{Скаларни свойства и именуване на контейнерите}

Реализацията на \code{property<T, Name>} свързва стойностен тип с compile-time име. Функцията \code{column\_name()} дава еднозначен достъп до символното име на полето и се използва от конекторите при рендиране към SQL колони, document fields или други backend-специфични понятия. Тази конструкция е едновременно минималистична и достатъчна, защото не въвежда външна регистрация на метаданни.

Именуването на целевия контейнер за съхранение се реализира чрез \code{table\_name\_trait<T>}. Макар името да е наследено от ORM традицията, архитектурно то играе по-обща роля: в SQLite и PostgreSQL означава таблица, в MongoDB --- collection, в Redis --- key prefix, а в Neo4j --- label. По този начин библиотеката използва единен логически интерфейс за адресиране на физически различни структури.

Практически това означава, че connector авторът не извлича информация от външна схема или от runtime registry, а от типовата система. Това намалява количеството boilerplate и прави грешките в именуването по-лесни за локализиране в компилационната фаза или в unit тестовете.

\subsection{Field wrappers, правила и placeholder-и}

Полето в query израз се представя чрез \code{field<\&T::m>}. Това е \code{constexpr} wrapper над member pointer, който носи достатъчно информация за типа на таблицата, типа на стойността и името на колоната. Операторите върху \code{field} изграждат типизирани \code{Rule} обекти. Така изрази от вида \code{field<\&User::id> == Placeholder<int>\{\}} не са низове, а compile-time описания на предикати.

Placeholder механизмът има две форми. \code{Placeholder<T>} моделира анонимни параметри, които се консумират позиционно от \code{execute()} аргументите. \code{orm::ph<T, std::placeholders::\_N>} моделира индексирани параметри, които позволяват една runtime стойност да се използва на повече от една позиция. Това е особено важно при backend-и като SQLite и PostgreSQL, където има естествена поддръжка на индексно адресирани placeholder-и.

Тестовете \code{IndexedPlaceholderRendersQuestionMarkN}, \code{IndexedPlaceholderReusedSameArgInSQL} и \code{TwoDistinctIndexedPlaceholdersRenderCorrectSlots} в \path{tests/integration/test_mockdb.cpp} показват, че тази подсистема не е теоретична добавка, а реално използван механизъм за изразяване на по-точни parameter binding сценарии.

\subsection{Изграждане на query обекти}

Фабричните функции \code{select}, \code{insert}, \code{update} и \code{deleteq} създават query обекти, върху които последователно се наслагват \code{where}, \code{join}, \code{order\_by}, \code{group\_by} и \code{limit}. Съществено е, че всяка от тези операции не променя низ, а връща нов типизиран обект с разширено вътрешно състояние. В резултат query builder-ът е едновременно fluent на вид и статичен по природа.

Това позволява в тестовете да се проверява не само крайният резултат от рендирането, а и самата форма на междинното представяне --- например броят на \code{where} клаузите или типът на response tuple-а. В \path{tests/integration/test_mockdb.cpp} това се вижда в \code{SelectWithInnerJoin}, \code{SelectWithOrderByDesc}, \code{SelectWithLimit} и множество други сценарии, които валидират различни фази на query composition.

Ключен имплементационен детайл е, че всички builder извиквания са \emph{константни изрази}. Вътрешният помощник \code{orm\_tuple::tuple\_cat} използва разгъване на пакет чрез индексни последователности вместо \code{std::tuple\_cat}, което прави целия builder chain constexpr-съвместим на всички поддържани компилатори, включително MSVC. Практическата последица е, че всеки query обект може да бъде деклариран \code{static constexpr}:

\begin{lstlisting}[language=C++,caption={Query обект деклариран като compile-time константа},label={lst:constexpr_query}]
using namespace orm::literals;
static constexpr auto get_active =
    orm::select(orm::field<&User::id>, orm::field<&User::name>)
        .where(orm::field<&User::score> > 0.0)
        .where(orm::field<&User::id> == orm::Placeholder<int>{})
        .limit(10_per_page & 2_page);
\end{lstlisting}

Операторът \code{\&} комбинира помощниците \code{\_per\_page} и \code{\_page} за user-defined literals в произволен ред и произвежда \code{Pagification} обект, чиято структура е известна по време на компилация. Декларирането на заявка като \code{static constexpr} я превръща в константа с нулева цена за изпълнение: IR на заявката се конструира веднъж и никога не се преизгражда при многократно изпълнение.

\begin{figure}[H]
\centering
\begin{tikzpicture}[node distance=1.35cm and 1.2cm]
    \node[ormaccent] (api) {\textbf{Fluent API}\\\code{select}/\code{where}/\code{join}/\code{limit}};
    \node[ormblock, below=of api] (typed) {Ново типизирано\\query състояние при всяка стъпка};
    \node[ormblock, below left=of typed] (shape) {Compile-time\\структура на заявката};
    \node[ormblock, below right=of typed] (params) {Runtime\\стойности на параметрите};
    \node[ormblock, below=of $(shape)!0.5!(params)$] (exec) {\code{connector\_trait<DB>::execute}};

    \draw[ormline] (api) -- (typed);
    \draw[ormline] (typed) -- (shape);
    \draw[ormline] (typed) -- (params);
    \draw[ormline] (shape) -- (exec);
    \draw[ormline] (params) -- (exec);
\end{tikzpicture}
\caption{Взаимодействие между query composition, compile-time структура и runtime параметри}
\label{fig:query_composition_execution}
\end{figure}

Фигура~\ref{fig:query_composition_execution} показва, че реализацията работи в два режима едновременно: структурната форма на заявката е фиксирана в типа, докато конкретните стойности се подават отделно. Точно това прави възможно едновременното наличие на compile-time проверки и runtime параметризация.

\subsection{Изпълнение, резултати и достъп до полета}

На етап изпълнение \code{db<DB>} предава query IR-а към \code{connector\_trait<DB>::execute}. Резултатът се връща като \code{orm::result<Row, Fields>}, където \code{Row} е tuple тип, а \code{fields} съдържа информация за избраните полета. Това решение позволява както итериране върху материализираните редове, така и достъп чрез \code{get\_field<\&T::m>}. Реализацията така остава типобезопасна и след изпълнението на заявката.

Тестовете \code{SelectResultRowTypeIsTuple}, \code{SelectMultiFieldRowType}, \code{SelectGetByMemberPointer} и \code{SelectGetByIndex} показват, че result слойът не е просто контейнер от анонимни редове, а структурирана type-safe абстракция. Това е особено важно за тезата, защото демонстрира, че типовата дисциплина не приключва с query builder-а, а продължава и в result handling-а.

\subsection{Prepared queries като повторно използваем execution артефакт}

\code{prepared\_query<DB, Query>} разширява изпълнителния модел. Вместо да кешира SQL низ, той съхранява връзка към backend обекта и самия query IR. Това е естествено продължение на compile-time подхода: библиотеката третира query обекта като първичен артефакт, а не като временна стъпка по пътя към низ.

\begin{lstlisting}[language=C++,caption={Извадка от \code{prepared\_query<DB, Query>}},label={lst:prepared_query}]
template <typename DB, typename Query>
class prepared_query
{
public:
    prepared_query(DB& conn, Query q) : conn_(conn), query_(std::move(q)) {}

    template <typename... Params>
    auto execute(Params&&... params) const
    {
        return connector_trait<DB>::execute(
            conn_, query_, std::forward<Params>(params)...);
    }

    [[nodiscard]] const Query& query() const noexcept { return query_; }
};
\end{lstlisting}

Listing~\ref{lst:prepared_query} показва важен инженерeн избор: повторната употреба е организирана около IR-а, а не около текстовия резултат от рендирането. Това позволява една и съща логическа заявка да бъде изпълнявана многократно с различни параметри, без да се реконструира builder формата.

\begin{figure}[H]
\centering
\begin{tikzpicture}[node distance=1.35cm and 1.2cm]
    \node[ormaccent] (build) {Създаване на query IR};
    \node[ormblock, right=of build] (prepare) {\code{db.prepare(query)}};
    \node[ormblock, right=of prepare] (store) {Съхранен IR +\\референция към \code{DB}};
    \node[ormblock, below=of store] (exec1) {\code{execute(1)}};
    \node[ormblock, left=of exec1] (exec2) {\code{execute(99)}};
    \node[ormblock, right=of exec1] (exec3) {\code{execute(7)}};

    \draw[ormline] (build) -- (prepare);
    \draw[ormline] (prepare) -- (store);
    \draw[ormline] (store) -- (exec1);
    \draw[ormline] (store) -- (exec2);
    \draw[ormline] (store) -- (exec3);
\end{tikzpicture}
\caption{Lifecycle на \code{prepared\_query}: един IR, множество изпълнения}
\label{fig:prepared_query_lifecycle}
\end{figure}

Точно това поведение се валидира от \code{PrepareReturnsExecutableObject}, \code{PrepareWithParamExecutesCorrectly}, \code{PrepareCanBeCalledMultipleTimesWithDifferentParams} и \code{PrepareExposesQueryIR} в \path{tests/integration/test_mockdb.cpp}. Тестовете доказват, че prepared query обектът е едновременно изпълним и inspectable, което е силно доказателство за консистентността на подсистемата.

\subsection{Миграционен слой и DDL генериране}

Реализацията на \code{migrate<DB>} сравнява описаната от entity типовете схема с \code{live\_schema}. Когато има разлика, се създават DDL операции и се подават към \code{connector\_trait<DB>::ddl\_for(...)}. Това е добър пример за clean separation на responsibilities: diff логиката е обща, а конкретният DDL синтаксис е backend-специфичен.

\begin{lstlisting}[language=C++,caption={Извадка от миграционния diff pipeline},label={lst:migration_diff}]
[[nodiscard]] std::vector<ddl_op> diff(
    const live_schema& live,
    std::span<const entity_meta> entities) const
{
    std::vector<ddl_op> ops;

    for (const auto& entity : entities)
    {
        auto live_it = live.find(entity.table_name);
        if (live_it == live.end())
        {
            ops.push_back(create_table_op{entity.table_name, entity.columns});
            continue;
        }

        for (const auto& [col, type] : entity.columns)
        {
            if (live_it->second.find(col) == live_it->second.end())
                ops.push_back(add_column_op{entity.table_name, col, type});
        }
    }
    return ops;
}
\end{lstlisting}

Интересно е, че именно тук се вижда разликата между \emph{минимален работещ конектор} и \emph{пълен конектор}. Един backend може да изпълнява заявки, без да поддържа DDL генериране. Но ако трябва да участва в schema-aware tooling, неговият contract трябва да бъде разширен.

\begin{figure}[H]
\centering
\begin{tikzpicture}[node distance=1.35cm and 1.2cm]
    \node[ormaccent] (entity) {Entity metadata\\\code{entity\_meta}};
    \node[ormblock, right=of entity] (live) {Runtime schema\\\code{live\_schema}};
    \node[ormblock, below=of $(entity)!0.5!(live)$] (diff) {\code{migrate<DB>::diff}};
    \node[ormblock, below left=of diff] (ops) {DDL операции\\\code{create/add/drop/alter}};
    \node[ormblock, below right=of diff] (ddl) {\code{generate\_ddl}};
    \node[ormblock, below=of $(ops)!0.5!(ddl)$] (run) {dry-run или execution path};

    \draw[ormline] (entity) -- (diff);
    \draw[ormline] (live) -- (diff);
    \draw[ormline] (diff) -- (ops);
    \draw[ormline] (diff) -- (ddl);
    \draw[ormline] (ops) -- (run);
    \draw[ormline] (ddl) -- (run);
\end{tikzpicture}
\caption{Поток на schema migration подсистемата}
\label{fig:migration_pipeline}
\end{figure}

Тестовете в \path{tests/unit/test_schema_migration.cpp} покриват създаване на таблица, добавяне и отпадане на колони, dry-run кодове на завършване, генериране на DDL и state machine поведението на \code{run()}. Това е добър пример как library subsystems с по-високо ниво на отговорност се валидират чрез насочени unit тестове, а не само чрез интеграционен smoke test.

\subsection{Нишкова безопасност и transaction helpers}

В библиотеката вече съществува отделен слой за \code{connection\_pool<DB, N>}, \code{thread\_local\_db<DB>} и \code{transaction\_guard<DB>}. Тези компоненти показват, че реализацията не спира до еднократно изпълнение на заявка, а разглежда и по-реалистични operational сценарии. Важното архитектурно решение е, че достъпът до тези механизми също е capability-базиран. Например \code{connection\_pool} изисква \code{supports\_concurrent\_execute}, а \code{transaction\_guard} --- \code{supports\_transactions}.

\begin{lstlisting}[language=C++,caption={Извадка от thread-safety helper слоя},label={lst:thread_safety_helpers}]
template <typename DB, std::size_t N>
class connection_pool
{
    static_assert(
     requires { typename connector_trait<DB>::supports_concurrent_execute; },
     "connection_pool requires connector_trait<DB>::supports_concurrent_execute.");

public:
    [[nodiscard]] connection_guard<DB> acquire();
};

template <typename DB>
class transaction_guard
{
    static_assert(
     requires { typename connector_trait<DB>::supports_transactions; },
     "transaction_guard requires connector_trait<DB>::supports_transactions.");
};
\end{lstlisting}

Така compile-time контрактът не се ограничава до query формата, а достига и до lifecycle политиките на конектора. Това е важно за thesis аргумента, че библиотеката не е само builder за заявки, а инфраструктура за систематична и безопасна работа с backend ресурси.

\begin{table}[H]
\centering
\small
\caption{Основни подсистеми и тяхната тестова опора}
\label{tab:core_subsystem_tests}
\begin{tabularx}{\textwidth}{L{3.1cm}L{3.8cm}Y}
\toprule
\textbf{Подсистема} & \textbf{Основни файлове} & \textbf{Тестова опора} \\
\midrule
Entity моделиране & \path{ORM/entity/property.hpp}, \path{ORM/entity/relationship.hpp} & \path{tests/unit/test_property.cpp}, \path{tests/unit/test_relationship.cpp} \\
Query composition & \path{ORM/query/select.hpp} и свързаните query headers & \path{tests/unit/test_query.cpp}, \path{tests/integration/test_mockdb.cpp} \\
Prepared queries & \path{ORM/connector/prepared_query.hpp} & \code{Prepare*} сценариите в \path{tests/integration/test_mockdb.cpp} \\
Schema migration & \path{ORM/db/migration/migration.hpp} & \path{tests/unit/test_schema_migration.cpp} \\
Thread safety & \path{ThreadSafety/thread_safety.hpp} & \path{tests/unit/test_thread_safety.cpp} \\
\bottomrule
\end{tabularx}
\end{table}

\subsection{Извод за реализацията}

Основните подсистеми имат пряка опора в кода и тестовете. \path{tests/unit/test_property.cpp} и \path{tests/unit/test_relationship.cpp} валидират entity описанията. \path{tests/unit/test_query.cpp} доказва, че query builder-ът наистина конструира правилни compile-time структури. \path{tests/integration/test_mockdb.cpp} и \path{tests/integration/test_sqlite.cpp} валидират изпълнението и prepared query сценариите. \path{tests/unit/test_schema_migration.cpp} покрива diff логиката на миграциите, а \path{tests/unit/test_thread_safety.cpp} проверява нишковия слой.

Следователно реализацията на основните подсистеми не е само проектно намерение, а архитектурно решение, подкрепено от систематична верификация. Това е и най-важният извод от настоящата глава: compile-time ORM архитектурата в хранилището не остава на равнище абстракция, а се материализира в работещи подсистеми с ясно разграничени отговорности и тестово покритие.
